\documentclass[track=junior,level=3]{kidscode}
%% one Scratch block in a table cell (stacked blocks: \sbk{...}\par\sbk{...})
\newcommand{\sbk}[1]{\begin{scratch}[scale=0.7]#1\end{scratch}}
%% the twin table: Scratch, Python, Scratch, Python
\newcolumntype{S}{>{\RaggedRight}m{\dimexpr0.24\linewidth-1.5\tabcolsep}}
\newcolumntype{P}{>{\RaggedRight}m{\dimexpr0.26\linewidth-1.5\tabcolsep}}
%% Python code in a cell: \pyl{t.forward(100)}  (an indented line: \pyin{...})
\newcommand{\pyl}[1]{{\ttfamily\bfseries #1}\par}
\newcommand{\pyin}[1]{\hspace*{2.2em}{\ttfamily\bfseries #1}\par}
\newcommand{\sheethead}[1]{\par\smallskip\noindent{\large\bfseries\color{kcaccent}#1}\par\vspace{2pt}}
%% the dyslexia edition keeps its font with slightly tighter lines so the sheet stays short
\makeatletter\ifkc@dyslexia\AtBeginDocument{\setstretch{1.15}}\fi\makeatother
\begin{document}
\kctitle{Cheat Sheet}{Scratch and Python side by side, running, and fixing bugs}
\vspace{-1.2cm}

\sheethead{Every block has a Python twin}
{\renewcommand{\arraystretch}{1.15}\small
\noindent\begin{tabular}{@{}SPS>{\RaggedRight\arraybackslash}m{\dimexpr0.26\linewidth-1.5\tabcolsep}@{}}
\toprule
\textbf{\faPuzzlePiece\ Scratch} & \textbf{\faCode\ Python} & \textbf{\faPuzzlePiece\ Scratch} & \textbf{\faCode\ Python}\\ \midrule
\sbk{\blockmove{move \ovalnum{100} steps}} & \pyl{t.forward(100)} &
  \sbk{\blockmove{turn \turnleft{} \ovalnum{90} degrees}}\par\sbk{\blockmove{turn \turnright{} \ovalnum{90} degrees}} & \pyl{t.left(90)}\pyl{t.right(90)}\\ \midrule
\sbk{\blockpen{set pen color to \pencolor{kcblue}}} & \pyl{t.color("blue")} &
  \sbk{\blockpen{set pen size to \ovalnum{5}}} & \pyl{t.pensize(5)}\\ \midrule
\sbk{\blockpen{pen up}}\par\sbk{\blockpen{pen down}} & \pyl{t.penup()}\pyl{t.pendown()} &
  \sbk{\blockmoreblocks{define \ovalmoreblocks{flower}}}\par(a My Block) & \pyl{def flower():}(stretch)\\ \midrule
\multicolumn{2}{@{}>{\RaggedRight}m{\dimexpr0.5\linewidth-\tabcolsep}}{\sbk{\blockrepeat{repeat \ovalnum{4}}{\blockmove{move \ovalnum{100} steps}\blockmove{turn \turnleft{} \ovalnum{90} degrees}}}} &
  \multicolumn{2}{>{\RaggedRight\arraybackslash}m{\dimexpr0.5\linewidth-\tabcolsep}@{}}{\normalsize\pyl{for i in range(4):}\pyin{t.forward(100)}\pyin{t.left(90)}%
   \smallskip{\small The lines inside the loop are \textbf{pushed in} 4 spaces. Press \keys{Tab}.\par}}\\
\bottomrule
\end{tabular}\par}

\noindent\begin{minipage}[t]{0.53\linewidth}
\sheethead{Run a program in Thonny}
\begin{enumerate}[leftmargin=1.8em,itemsep=1pt,topsep=0pt]
  \item[\faFolderOpen] Click \menu{Load}. Open \texttt{starter.py}.
  \item[{\color{kcgreen}\faPlayCircle}] Click the green \menu{Run}.
  \item[\faEye] Look at the turtle window.
  \item[\faTimesCircle] Close it before you Run again.
  \item[\faSave] Click \menu{Save} when you finish.
\end{enumerate}
{\small Run is the green flag of Python. No window? Press \keys{Cmd}+\keys{Tab}.\par}
\end{minipage}\hfill
\begin{minipage}[t]{0.43\linewidth}
\sheethead{Fixing a bug}
\begin{enumerate}[leftmargin=1.6em,itemsep=1pt,topsep=0pt]
  \item[\faPlay] \textbf{Test.} Run it.
  \item[\faEye] \textbf{Look.} What did it do?
  \item[\faWrench] \textbf{Change one thing.}
  \item[\faRedo] \textbf{Test again.}
\end{enumerate}
{\small Bugs are clues. Everyone gets them!\par}
\end{minipage}

\sheethead{Red text is a clue}
Read the \textbf{last} red line. It says the line number.
\enquote{Did you mean} tells you the right word.
The Helper reads the line letter by letter. Check the line above it too.\par\smallskip
{\renewcommand{\arraystretch}{1.2}\small
\begin{tabularx}{\linewidth}{@{}>{\RaggedRight}p{0.42\linewidth}>{\RaggedRight\arraybackslash}X@{}}
\toprule
\textbf{Bug} & \textbf{Fix}\\ \midrule
\texttt{t.forwad(100)} & Spell it: \texttt{t.forward(100)}\\
\texttt{t.left(90} & Add the missing \texttt{)}\\
\texttt{t.color("bleu")} & Spell the color: \texttt{"blue"}\\
A line is not pushed in & Press \keys{Tab} at the start of the line\\
\bottomrule
\end{tabularx}\par}

\noindent\begin{minipage}[t]{0.56\linewidth}
\sheethead{The turn card}
\textbf{quarter turn = 90}\quad\textbf{half turn = 180}\par\smallskip
%% alt: The turn card: a triangle labeled 120, a square labeled 90 and a hexagon labeled 60.
\begin{tikzpicture}[font=\bfseries]
  \node[regular polygon,regular polygon sides=3,draw=black,fill=kcorange!40,minimum size=1.4cm] (t) at (0,0) {};
  \node[below=1pt of t] {120};
  \node[regular polygon,regular polygon sides=4,draw=black,fill=kcsky!40,minimum size=1.4cm] (s) at (2.0,0) {};
  \node[below=1pt of s] {90};
  \node[regular polygon,regular polygon sides=6,draw=black,fill=kcgreen!40,minimum size=1.4cm] (h) at (4.0,0) {};
  \node[below=1pt of h] {60};
  \node[font=\small] at (0,1.0) {triangle};
  \node[font=\small] at (2.0,1.0) {square};
  \node[font=\small] at (4.0,1.0) {hexagon};
\end{tikzpicture}
\end{minipage}\hfill
\begin{minipage}[t]{0.40\linewidth}
\sheethead{Colors I can type}
{\renewcommand{\arraystretch}{1.1}\begin{tabular}{@{}l@{\hspace{1.2em}}l@{}}
\swatch{kcorange}\ \texttt{"orange"} & \swatch{kcblue}\ \texttt{"blue"}\\
\swatch{kcpurple}\ \texttt{"purple"} & \swatch{kcgreen}\ \texttt{"green"}\\
\swatch{kcyellow}\ \texttt{"gold"} & \swatch{kcblack}\ \texttt{"black"}\\
\end{tabular}\par}
{\small Small letters, with quote marks on both sides.\par}
\end{minipage}

\end{document}
